\section{Which Evidence Supports Defense Transfer?}\label{sec:evaluation}\label{sec:threats}
\subsection{Separate the evaluated object before counting holdouts}
RQ4 distinguishes three questions: does an \emph{attack} recognize new inputs or systems; is a \emph{defense condition} evaluated; and does a \emph{fixed attacker or defense calibration} carry from a specified source to a target? Table~\ref{tab:r26-transfer-objects} makes the distinction operational. Repeating a defense on several models can establish condition-specific outcomes without establishing a held-out model or transferring a fitted calibration.
\begin{table}[t]\centering\footnotesize
\caption{Three evidence objects for RQ3. The same family may supply several objects in separate experiments.}\label{tab:r26-transfer-objects}
\begin{tabularx}{\columnwidth}{@{}p{.24\columnwidth}YY@{}}\toprule
Object & Required link & Permitted reuse\\\midrule
Attack generalization & Training/profiling source $\to$ held-out inputs or systems; attack outcome. & An attacker or recognition method covers the stated target domain.\\
Defended condition & Named mechanism, local fitting, defended test and outcome. & Retain that operating point; multi-model testing alone is not held-out transfer.\\
Defense transfer & Source-selected policy/calibration held fixed $\to$ target; local-calibration comparator and attacker fit specified. & Transport the evaluated policy only over the tested direction and scope.\\\bottomrule
\end{tabularx}\end{table}


The 82-family matrix in Appendix~\ref{app:holdouts} is a discovery index for these questions. Its six axes are task/input/user, model, harness/topology, hardware, network/load and time. H marks reported separation, V condition variation/recalibration, X a scoped restriction, and ? unresolved extraction. The 13 task/input H entries include only four mitigation families; all three model H entries concern attack/identification generalization in families outside the 20-family mitigation audit. The two topology H entries occur in mitigation families, but their held-out objects are workflow recognizers. These counts therefore cannot be read as 13 task-transfer or three model-transfer validations of defenses. S34 identifies each H entry and its evidence object; S28 preserves the full matrix and source scopes.

\subsection{A changed generator and a tested defense are different bindings}
LLM-Fin, ArchSnoop and LeakyLMs evaluate recognition of unseen model variants or architectures \cite{nazari2024llmfin,weng2026archsnoop,majidi2026leakylms}. This supplies attack-side model generalization. In the traffic benchmark, graph-to-star and star-to-graph classification explicitly move the fitted attacker across topologies; the prompt/temperature countermeasures are separate evaluations \cite[V-B--C]{pouryousef2026traffic}. Dangol's leave-composition-out test likewise examines recognition of a novel workflow shape, separately from its protected-capture frontier \cite[9.8--9.12]{dangol2026workflow}. Neither topology entry licenses a transported defense calibration.

There is nevertheless useful \emph{defense-side variation} evidence. Time Will Tell applies the same fixed-length instruction to two models and reports different residual classification ASRs; this tests a common rule under changed generators \cite[VII-B]{timewilltell2024}. Whisper Leak evaluates mitigation conditions across provider-models, while its prompt holdout describes the classifier protocol \cite[3.2,6]{whisperleak2025}. Preserve those per-condition results; a shared rule, a refitted local attack, a source-trained attack and a source-calibrated defense are four distinct comparisons.

\paragraph{Deployment decision}
For the proposed change, name the source, target, held-out unit, quantities held fixed or refitted, and protected outcome. Carry forward only the corresponding binding. An attack-generalization result justifies including that attacker at the target; a defended-condition result supplies its local outcome; a source-to-target defense-calibration result supports only the stated transfer. This prevents both overclaiming portability and discarding useful local evidence. Unresolved hardware/time entries in the discovery matrix identify source follow-up, not absent experiments.
